% =============================================================================
% Abstract
% =============================================================================

\begin{abstract}
\setlength{\parindent}{0pt}

\noindent Diese praxisorientierte Fallstudie untersucht die Verfügbarkeit und Fehlertoleranz des ROSI-Systems (\textit{Realtime Optical Surface Inspection}) im industriellen 24/7-Betrieb. ROSI ist ein on-premise Inline-Scan-System zur automatisierten Qualitätskontrolle von Baustoffen: Materialoberflächen werden erfasst, Defekte mittels KI-Inferenz erkannt und anschließend regelbasiert in Qualitätsstufen eingestuft.

\noindent Der Schwerpunkt liegt auf Fehlern, die im Dauerbetrieb schleichend entstehen, lange unbemerkt bleiben oder kaskadieren und im Entwicklungsbetrieb daher häufig nicht sichtbar werden. Methodisch werden fünf definierte Failure Cases auf Prozess- und Betriebssystemebene in einer kontrollierten Testumgebung reproduzierbar simuliert. Pro Case werden allgemeine Messgrößen erhoben, darunter Time-to-Detection bei Prozessausfällen, Inspektionsdurchsatz (Inspections/s) unter variierender Last sowie Queue-Tiefe und Speicherfüllstand über die Zeit.

\end{abstract}
